\section{Trabalhos Relacionados}
\label{sec:relacionados}

A literatura em arquitetura de computadores e HPC destaca a importância de avaliar sistemas de forma integrada, considerando processadores, memória, interconexões, armazenamento e aceleradores~\cite{hennessy2019golden}. Essa visão é útil para este artigo porque o desempenho observado em uma carga matricial não depende apenas do processador, mas também das bibliotecas numéricas e da forma como os dados são movimentados durante a execução.

Trabalhos sobre benchmarks de HPC, como avaliações baseadas em LINPACK, mostram a relevância de métricas padronizadas para comparar sistemas~\cite{dongarra2003linpack}. Embora LINPACK não represente sozinho todas as cargas modernas, ele permanece relevante por medir desempenho em operações de álgebra linear, uma família de operações também presente em redes neurais.

Um trabalho diretamente relacionado ao experimento deste artigo é o de Goto e van de Geijn~\cite{goto2008anatomy}, que discute como multiplicações de matrizes de alto desempenho dependem de bloqueio, uso eficiente de cache e movimentação de dados. Essa discussão ajuda a interpretar por que matrizes maiores podem apresentar melhor aproveitamento das rotinas numéricas: parte do desempenho vem menos da operação matemática isolada e mais da forma como a implementação organiza dados e instruções na arquitetura disponível.

Huerta et al.~\cite{huerta2020convergence} discutem a convergência entre IA e HPC, argumentando que plataformas de alto desempenho podem reduzir o tempo de experimentação e apoiar aplicações científicas baseadas em dados. Essa perspectiva é relevante para este trabalho por tratar a IA não apenas como algoritmo, mas como carga computacional que depende de infraestrutura adequada.

Outra linha de pesquisa analisa aceleradores especializados para aprendizado de máquina. Jouppi et al.~\cite{jouppi2017tpu} apresentam uma avaliação de unidade de processamento tensorial em data centers, evidenciando como arquiteturas dedicadas podem alterar significativamente o desempenho de cargas de inferência. Embora o presente estudo seja realizado em uma máquina local, essa discussão ajuda a contextualizar por que GPUs e aceleradores são componentes centrais em ambientes de IA generativa.

No campo de avaliação de desempenho, a suíte MLPerf Inference foi proposta para medir a rapidez com que sistemas executam modelos em diferentes cenários de implantação~\cite{reddi2019mlperf}. A documentação atual da suíte inclui cargas de linguagem de grande porte, como Llama 2 70B e Llama 3.1 405B, reforçando a relevância de benchmarks padronizados para comparar infraestruturas destinadas a IA generativa~\cite{mlperfInferenceDocs}.

Este artigo se diferencia por adotar uma abordagem experimental de pequena escala. O foco não é competir com avaliações de data centers ou clusters, mas mostrar como uma medição simples, feita em uma máquina local e documentada em repositório, pode apoiar a compreensão de conceitos de HPC aplicados a cargas numéricas associadas à IA generativa.
